\question[4]
Das folgende Programm soll aus \texttt{'Habicht'} den String \texttt{'KABICHT'} machen.

\begin{lstlisting}
s = 'Habicht'
s[0] = 'K'
s.upper()
print(s)
\end{lstlisting}

\begin{parts}
\part Zeile 2 führt zu einem Fehler. Erläutere, warum.
\part Auch ohne Zeile 2 würde \texttt{Habicht} ausgegeben und nicht \texttt{HABICHT}. Erläutere, warum.
\part Gib ein korrigiertes Programm an.
\end{parts}

\begin{solutionbox}{6cm}
\begin{parts}
\part Strings sind unveränderlich. Man kann einzelne Zeichen nicht durch eine Zuweisung ändern (\texttt{TypeError}).
\part \texttt{upper()} verändert den String nicht, sondern liefert einen neuen String. Dieser wird nicht gespeichert.
\part
\begin{lstlisting}
s = 'Habicht'
s = 'K' + s[1:]
s = s.upper()
print(s)
\end{lstlisting}
\end{parts}
\end{solutionbox}
